The foregoing developments were intended to demonstrate the vital importance of the rule of progression in a society whose number of participants increases considerably, and thus the truth of the assertion I made on p. [127]. To illuminate this subject from every side, however, it will be necessary to follow the history of the origin of that article of the statute step by step. For this purpose I have carefully read the relevant records, and read them repeatedly, and I give an extract from them insofar as they concern that article. Here and there I shall also have to bring out some incidental circumstances, subordinate in themselves, if they can help to make the course of these deliberations more intelligible. I shall first remark that in 1794 the annual contribution was 5 Thl. in gold, the widows' pension 110 Thl. Kassenmünze, and that orphans who had lost both father and mother were entitled to the pension until they had completed their \(12^{\text{th}}\) year.

In a ministerial rescript addressed to the University and dated 30 June 1794, reference was made to a gift of 1000 Thl.\footnote{Kassenmünze. According to the accounts for 1793, it was paid in on 18 April 1793.} which the King had made to the widows' fund some time previously, and notice was given of the grant of a second gift of 500 Thl. in gold. It was added that, since these gifts were intended to bring about an early increase in the pensions, the participants were expected to contribute on their part to attaining that end and to be willing to raise their annual contributions from 5 Thl. in gold to 10 Thl. Kassenmünze. According to the calculations in an attached P. M., there was no doubt that it was quite practicable to introduce an increase in the pensions immediately, so that the six oldest widows would henceforth receive, instead of the previous 110 Thl., 150 Thl., and all the others 130 Thl. each. In conclusion, an offer was made, should the widows' fund's capital not all be invested with complete security, to arrange its secure investment at 3 percent with public or municipal treasuries. A similar offer had already been made once before, in the comments on the annual accounts for 1792.

It is unnecessary for my purpose to go through the attached enclosure, whose author is not named, in detail. But I shall bring out a few incidental circumstances.

I. The widows' fund's capital, it says, was lent at unequal interest rates, some of it\footnote{At that time nearly the third part of the capital assets.} at 3 percent, other portions at higher rates. But since the interest rate on all the capital might easily fall to 3 percent, and estimates ought to reckon with sums as secure as possible, no more than 3 percent was to be assumed in the calculations.

This, together with the conclusion of the rescript just cited, explains why the University's deliberations likewise reckoned, without further remark, on only 3 percent for all future increases in capital; interest at \(3\tfrac12\) percent was entered only for the capital already present and already invested. Compare the note to p. [131.]

II. To estimate the number of widows on which the calculation should be based, the writer continues:

‘If, on the basis of previous experience, one now assumes that 3 existing marriages have to sup-\par
\noindent ‘port one widow, the 26 married professors (out of a total of 36) would have about 9 widows\par
\noindent ‘to maintain. But since it gives greater certainty to take the extreme case as the basis,\par
\noindent ‘one should rather reckon \(2\tfrac12\) marriages against \textit{one} widow, so that the existing\par
\noindent ‘26 professors' marriages would have to maintain — 10 widows.’
\par\SourceLeaf{138}
In the estimates made on the University's side, those 3 or \(2\tfrac12\) marriages to one widow were considered too many, and the ratio of 2 marriages to one widow was taken as the basis.

I have cited the passage from the P. M. here solely because it explains why people were so easily satisfied with the ratio of two marriages to one widow, although, very probably, even this is not yet suited to the circumstances of the professors' widows' fund, and still fewer marriages should be reckoned against one widow. More detail on the matter itself will follow below. But the course of the business involuntarily recalls (if I may draw an analogy here from a humbler sphere) buyers who, with an imperfect knowledge of the goods themselves, believe they have made a good bargain if they have bought far below the price first asked, although, on closer inspection, they still paid too much. I need not remind the reader that I do not wish this analogy to be stretched beyond its proper extent, for the unknown proposer undoubtedly put forward the ratios \(3:1\) and \(2\tfrac12:1\) in good faith as to their admissibility.

When the then Prorector F. circulated the rescript among the Senate on 6 July, he added a third subject for deliberation to the two it contained (increasing the contributions and increasing the pensions), by proposing to extend the duration of orphans' pensions until completion of the \(20^{\text{th}}\) year of life. He left it to the Senate members to express themselves on these matters immediately in writing, or at the Senate meeting fixed for 12 July.

This circular was signed by 17 Senate members; only a few of the statements made on it need to be mentioned here.

The then curator of the widows' fund, P., puts forward above all the principle that the fund owes present widows as much as future ones, but also owes future widows exactly as much as present ones. He explains this phrase (in itself indeed very vague) to mean that any future widow receiving less than another had received earlier would (in his opinion) truly suffer an injury, and that the first principle must therefore be to fix the pension amount so that, with the highest probability, it should never need to be reduced again. In this respect, however, P. considers the calculation in the rescript's enclosure insufficiently secure; one should reckon not \(2\tfrac12\) marriages to 1 widow but only 2, and should therefore set the maximum number of widows not at 10—12 but at 14—15, and consequently also assume smaller pension amounts.

\textsc{Kaestner} considers the question too involved and difficult for anything to be settled without a \textit{detailed} and \textit{precise} investigation; he too holds that only \textit{at most} 2 marriages to 1 widow may be reckoned. Since he had long since left the widows' fund (he had been a participant from 1755—1773), he had no vote in the matter (as senior of the Faculty of Philosophy he was nevertheless a member of the Church Deputation), but advised that no decision be taken without first consulting an expert, perhaps p. \textsc{Kritter}.

G. also wishes the fund's resources to be determined by precise calculations, and points to the \textit{terrible} consequences of hasty decisions to grant excessively large pensions, with the examples of the Hanoverian, Bremen and other widows' societies.

The remaining opinions partly agree with the preceding ones and partly develop misgivings about increasing the contributions or extending the orphans' pensions, which need not be excerpted here.

At the Senate meeting of 12 July, attended by 12 professors including the Prorector, 9 declared themselves unconditionally \textit{in favour of} increasing the contributions, 2 adding that they considered the increase too large; one (who had emphatically opposed it in the written vote) wished to join the majority. — F.'s proposal to extend the orphans' pensions was still deferred pending a closer consideration of the fund's circumstances. — As regards
\SourceLeaf{139}
increasing the widows' pensions, it was decided to obtain an opinion from \textsc{Kritter}\footnote{It does not appear that anything was decided as to \textit{the extent to which} \textsc{Kritter}'s advice was to be sought. \textsc{Kaestner} was not present at the meeting.}. Yet the meeting already declared itself \textit{against} the larger pension proposed in the rescript for the six oldest widows, since they were already favoured by V.'s bequest.

The consultation of p. \textsc{Kritter}, conducted through a letter from P. whose draft is among the records, was in fact a very limited one. After an account of the widows' fund's principal arrangements given only in very general outline, just two questions are put to \textsc{Kritter}: I) By what principle and ratio the maximum number of widows in a society of this kind should be determined, and II) by how much this maximum number of widows' pensions to be paid should be increased if, when there is no widow, or after her earlier death, entitlement to the pension passes to any orphans there may be and continues until the youngest child reaches the age of 20. The current number of all participants in the fund, and of the married men among them, is given, exactly as in the rescript's enclosure cited above, as 36 and 26. It is added that these numbers and their ratio are variable and could hardly be subjected to a rule of probability; the letter concludes with an offer to communicate immediately any further data the person consulted might need from previous experience concerning the ratio of participants to widows.

I enclose a complete copy of \textsc{Kritter}'s answer, or his ‘opinion’, dated 19 July\footnote{Opinion on some questions put to me concerning the Göttingen University widows' fund.

Ad I. The body of professors at Göttingen has existed for over 50 years since the University was founded, so that even 10 years ago it could have had the greatest number of widows who, according to the laws of mortality, must have existed for some 22 or 24 married professors, namely 12 widows. This has indeed been shown by the list communicated to me of the growth in the annually increasing number of widows, and would have been shown still better had the body of professors been able to consist of 100 persons. But since, by the nature of the matter, this number was only about the \(4^{\text{th}}\) part of 100, it was also natural that the number of widows from 1775 to 1780 exceeded 12, and from 1781 to 1792 was below 12, because with small numbers the order of mortality cannot occur as it does with large numbers. Nevertheless, one must still assume that the greatest number of widows, in an average of some 10 years, will constantly be about half the number of married professors.

Ad II. Since it is now desired that the orphaned families left by a deceased widow or also by a deceased widower should receive a pension equal to a widow's until completion of the \(20^{\text{th}}\) year, I can determine something about this only from the experience of the Bremen widows' fund. The Bremen society had to expect about the \(6^{\text{th}}\) part more in widows' pensions when it decided to pension orphaned families like a widow and to continue doing so until completion of the youngest child's \(18^{\text{th}}\) year. But since the orphaned families of the professors are to enjoy this until completion of the youngest child's \(20^{\text{th}}\) year, one must surely reckon, against 12 widows' pensions, more than 2 orphans' pensions.

I therefore consider it advisable to add to 12, as the maximum number of widows, at least 2, so that the divisor by which the annual income from the fund's assets and other sources is divided is set at 14. Since there are now only 10 widows, the whole annual revenue of the fund must not be divided among these 10, because otherwise the widows who join later would be shortchanged; instead one must divide by 14. Then within 10 years the fund will have been strengthened to about 25 widows' shares, the interest on which will suffice for one pension more than before, so that, when the true maximum of widows' and orphans' pensions occurs, it will be possible to pay 15 pensions; and should anything remain over, an addition could be given especially to poor widows.

Some allowance ought surely also to be made for possible misfortunes affecting the invested capital, loss of interest and other expenses.

Göttingen, 19 July 1794.\hfill J. A. \textsc{Kritter}.}.

